\documentclass[runningheads]{llncs}
\usepackage{graphicx} % Required for inserting images

\usepackage{amsmath,amssymb}
\usepackage{url}

\title{LSTA for clean and dirty ancilla}
\author{Yu-Fang Chen}
\authorrunning{Yu-Fang Chen}
\institute{Institute of Information Science, Academia Sinica, Taiwan}

\begin{document}

\maketitle

\section{Introduction}

\section{Preliminaries}

\subsection{Level-synchronized tree automata}

An $n$-qubit vector $\lvert\psi\rangle=\sum_{b\in\{0,1\}^n}
\alpha_b\lvert b\rangle$ is encoded by a perfect binary tree of depth $n$.
Internal nodes at depth $i-1$ carry $x_i$; the leaf reached by $b$ carries
$\alpha_b$. Left and right edges encode $0$ and $1$, respectively.

\begin{definition}[LSTA]
Following~\cite{lsta}, an LSTA is a tuple
$\mathcal A=(Q,\Sigma,\Delta,\mathcal R)$, with finite state set $Q$,
root states $\mathcal R\subseteq Q$, ranked alphabet
$\Sigma=\Sigma_2\uplus\Sigma_0$, and finite transition set $\Delta$.
Transitions have the forms
\[
 q\xrightarrow{C} f(q_0,q_1)
 \qquad\text{or}\qquad
 q\xrightarrow{C} a,
\]
where $f\in\Sigma_2$, $a\in\Sigma_0$, $q,q_0,q_1\in Q$, and
$C\subseteq\mathbb N_0$ is finite.
Distinct transitions with the same source state have disjoint choice sets.
\end{definition}

A run assigns a matching transition to every node, respecting the ordered
child states. It is accepting if its root state belongs to $\mathcal R$
and, at every depth $d$ (including leaves),
\[
 \bigcap_{\substack{v\text{ a node}\\ |v|=d}} C_{\rho(v)}
 \neq\varnothing.
\]
Thus all transitions at one level share a choice; different levels may
use different choices. The language $\mathcal L(\mathcal A)$ consists of
trees admitting accepting runs. For quantum states, we use qubit labels
internally and complex amplitudes at leaves, and restrict to normalized
vectors with the prescribed qubit order.

\subsection{Clean and dirty ancillas}

Let $D$ be an $n$-qubit data register, $A$ an $m$-qubit ancilla register,
and $U$ a unitary on $\mathcal H_D\otimes\mathcal H_A$.
An ancilla provides temporary workspace rather than part of the intended
output. We distinguish its initialization from the condition required for
safe reuse~\cite{gidney,dirty}.

\begin{definition}[Clean-ancilla safety]
The ancilla register $A$ starts in $\lvert0^m\rangle_A$ and must return
to $\lvert0^m\rangle_A$ for every data input.
Formally, $U$ satisfies clean-ancilla safety if there exists a unitary $V$
on $\mathcal H_D$ such that, for every state $\lvert\psi\rangle_D$,
\begin{equation}\label{eq:clean-safety}
 U(\lvert\psi\rangle_D\otimes\lvert0^m\rangle_A)
 =V\lvert\psi\rangle_D\otimes\lvert0^m\rangle_A.
\end{equation}
\end{definition}

\begin{definition}[Dirty-ancilla safety]
The ancilla register $A$ starts in an arbitrary unknown state
$\lvert\phi\rangle_A$ and must return to the same state.
Formally, $U$ satisfies dirty-ancilla safety if there exists a unitary $V$
on $\mathcal H_D$ such that, for every pair of states
$\lvert\psi\rangle_D$ and $\lvert\phi\rangle_A$,
\begin{equation}\label{eq:dirty-safety}
 U(\lvert\psi\rangle_D\otimes\lvert\phi\rangle_A)
 =V\lvert\psi\rangle_D\otimes\lvert\phi\rangle_A.
\end{equation}
The same $V$ must work for all inputs; in particular, it cannot depend on
$\lvert\phi\rangle_A$.
\end{definition}

By linearity, the dirty-ancilla condition also covers inputs in which
the ancilla is entangled with other registers, including an untouched
external system.

Restoring only the basis values of $A$ is insufficient:
$I_D\otimes Z_A$ leaves those values unchanged but maps
$\lvert+\rangle_A$ to $\lvert-\rangle_A$.
Dirty-ancilla safety implies clean-ancilla safety, whereas the converse
fails by this example. The existential $V$ absorbs any overall global phase.
If a target operation is specified, its correctness is a separate
requirement on $V$.

\begin{thebibliography}{9}
\bibitem{lsta}
Parosh Aziz Abdulla et al.
\emph{Verifying Quantum Circuits with Level-Synchronized Tree Automata
(Technical Report)}.
arXiv:2410.18540, 2024.
\url{https://arxiv.org/abs/2410.18540}.

\bibitem{gidney}
Craig Gidney.
\emph{Factoring with $n+2$ clean qubits and $n-1$ dirty qubits}.
arXiv:1706.07884, 2017.
\url{https://arxiv.org/abs/1706.07884}.

\bibitem{dirty}
Bonan Su, Li Zhou, Yuan Feng, and Mingsheng Ying.
\emph{Borrowing Dirty Qubits in Quantum Programs}.
arXiv:2508.17190, 2025.
\url{https://arxiv.org/abs/2508.17190}.
\end{thebibliography}

\end{document}

